\section{Signed memory and the determinant moment}
\label{mem:section}

We retain the notation, physical state space, good-state tests, and
operator $\mathcal A$ of the preceding section. In particular, the number
$J=\lfloor L^{.01}\rfloor$ of pads per group grows with $x$,
$M=J+1$, and $N=2R\asymp L^{.5}$. All constants depending on the fixed
number $K$ of groups are allowed to affect the threshold for $x$.
We will explicitly distinguish those constants from powers of $L$.

\begin{proposition}[The minor-arc moment]\label{prop:minor-moment}
For every fixed $E_0>0$, one can choose $A_0$ sufficiently large and then
$K$ sufficiently large, in terms of $E_0,q$ and the fixed parameters of
Theorem~\ref{thm:type-II}, so that the operator with the good-state
projections satisfies
\begin{equation}\label{mem:moment-bound}
 \frac{1}{UV}\sum_{\substack{P\in\Omega\\P\ \mathrm{primitive}}}
 \inner{\mathbf u_P}{(\mathcal A\mathcal A^*)^R\mathbf u_P}
 \le L^{-E_0N}.
\end{equation}
The choices of $A_0,K$ do not depend on the particular fixed band
exponents $a_1<\cdots<a_K$; the threshold for $x$ may depend on them.
Consequently, for any prescribed fixed $A>0$, the determinant indicator
in the expanded square of the preceding section can be replaced by
\eqref{eq:det-major}, with error $O_A(XY L^{-A})$. In the resulting
major term the disjointness restriction on the two unshared products
$a,b$ can be omitted. Equivalently, with the notation of
\eqref{det:cauchy-square} and \eqref{maj:sum},
\begin{equation}\label{mem:replacement-interface}
 \abs{\mathcal Q_Y-\mathcal Q_Y^{\mathrm{maj}}}
      \ll_A XY L^{-A}.
\end{equation}
The precision needed in the pairing estimate
\eqref{eq:det-pairing} for this consequence is independent of $K$.
\end{proposition}

The proof constructs operators which remember a prime between two of
its uses. The construction has a resemblance to the lifespan expansion
in \cite[Section~3]{SmoothShifted}; the primewise identity and every
operator estimate needed here are proved below for the determinant
graph.

\subsection{The exact primewise identity}

Apply Lemma~\ref{det:root-replacement} to the expansion
of the left side of \eqref{mem:moment-bound}. We work for now at a fixed
archimedean root and with independent uniform kernel lines at the group
primes. Write the path as $z_0,\ldots,z_N$, where $z_0=z_N=e_1$;
all the $z_j$ are primitive and satisfy the tube and box conditions of
\eqref{det:tube}. Put
\begin{equation}\label{mem:baseline}
 q_j=\begin{cases}\sqrt q,&j=0,N,\\q,&0<j<N,\end{cases}
 \qquad \eta'_j=1-q_j,\qquad
 b'_p=1-\frac{\sum_{j=0}^N\eta'_j}{p+1},\qquad
 \nu_i(p)=\frac{1}{V_i(p+1)b'_p}\quad(p\in\mathcal P_i).
\end{equation}
Thus $b'_p>0$ for large $x$. With
$p_{\min}=\min\bigcup_i\mathcal P_i$, we have
\begin{equation}\label{mem:measure-masses}
 \frac{\nu_i(p)}{\mu_i(p)}=1+O\!\left(\frac{N+1}{p}\right),\qquad
 s_i:=\sum_{p\in\mathcal P_i}\nu_i(p)
       =1+O\!\left(\frac{N+1}{p_{\min}}\right),\qquad
 \sup_{i,p}\nu_i(p)\le C_K e^{-L^{.1}}.
\end{equation}
The estimates are uniform in the path. Here and below $C_K$ can also
depend on $q$ and on fixed smooth cutoffs.

For one prime $p$, let $A_p$ be its set of active visits, namely the
visits where $p$ occurs in the active list. For a line
$\ell\in\mathbb P^1(\mathbb F_p)$ set
$H_j(\ell)=\1_{\ell=[z_j]_p}$. Before active-list normalizations, its
factor in the independent-line expectation is exactly
\begin{equation}\label{mem:prime-factor}
 \frac{1}{p+1}\sum_{\ell\in\mathbb P^1(\mathbb F_p)}
 \prod_{j\in A_p}H_j(\ell)
 \prod_{j\notin A_p}(1-\eta'_j H_j(\ell)).
\end{equation}
Indeed the physical damping charges $q_j$ precisely when the prime
divides the first coordinate at visit $j$ without belonging to its
active list.

First suppose $A_p=\varnothing$. Expand the product in
\eqref{mem:prime-factor}. The empty and singleton subsets contribute
$b'_p$. In every larger subset fix its first and last indices $a<c$;
the sum over all choices of internal indices gives
\begin{equation}\label{mem:orphan-identity}
 b'_p+\frac{1}{p+1}\sum_{0\le a<c\le N}
 \eta'_a\eta'_c\,\1_{[z_a]_p=[z_c]_p}
 \prod_{a<j<c}\bigl(1-\eta'_j\1_{[z_j]_p=[z_a]_p}\bigr).
\end{equation}
After division by $b'_p$, a summand is a lifespan beginning with a
ghost at $a$ and ending with a ghost at $c$. It has one factor
$((p+1)b'_p)^{-1}$, one factor $-\eta'_j$ at each endpoint, and a
factor $q_j$ at each strictly internal visit that hits its line.
There is no lifespan with just one ghost: all singleton terms have
already been included in the baseline.

Next suppose $A_p\ne\varnothing$. If its active lines disagree, the
factor is zero. Otherwise let $\ell$ be their common line and put
$a=\min A_p$, $c=\max A_p$. This line is chosen with probability
$1/(p+1)$. The inactive visits between $a$ and $c$ retain their damping
factors. The products outside this interval have the exact identities
\begin{align}
 \prod_{j<a}(1-\eta'_jH_j(\ell))
 &=1-\sum_{h<a}\eta'_h H_h(\ell)
                \prod_{h<j<a}(1-\eta'_jH_j(\ell)),\label{mem:prefix}\\
 \prod_{j>c}(1-\eta'_jH_j(\ell))
 &=1-\sum_{h>c}\eta'_h H_h(\ell)
                \prod_{c<j<h}(1-\eta'_jH_j(\ell)).\label{mem:suffix}
\end{align}
These follow by telescoping a finite product, in opposite orders at the
two ends. They attach either no ghost or one ghost extension to each
end of the active interval. Every inactive internal hit then pays
$q_j$, and each chosen ghost endpoint pays $-\eta'_j$.

Extract $\prod_p b'_p\le1$ over all group primes. Equations
\eqref{mem:orphan-identity}--\eqref{mem:suffix} show that every term
of the expansion assigns at most one lifespan to each prime. All
active visits and ghost endpoints of a lifespan must have the same
line. Its probability
$1/(p+1)$, together with division by $b'_p$, contributes the factor
$((p+1)b'_p)^{-1}$ exactly once. The remaining factors are the damping
at inactive internal visits and the active-list normalization.
An \emph{active run} is a maximal interval of consecutive active visits.
The ban in the physical edge ensures that a prime active at two
consecutive visits is continued in a shared slot. Thus a lifespan in
group $i$ with $r$ active runs has factor
\begin{equation}\label{mem:run-charge}
 \frac{V_i^{-r}}{(p+1)b'_p}
\end{equation}
before its ghost signs and internal damping. The power of $V_i^{-1}$
counts active runs; the line probability is paid only once.

A lifespan retains its prime and line while the prime is outside the
active list; call such a retained prime \emph{pending}. At a strictly
internal inactive visit it is a hit if $[z_j]_p$ equals that line, and
a miss otherwise. The prime keeps the same line across both kinds of
visit.
Figure~\ref{fig:lifespan} shows how one lifespan can join two active
runs across a pending hit and a pending miss.
\begin{figure}[tb]
\centering
\begin{tikzpicture}[x=1.64cm,y=1cm,
  hit/.style={circle,draw=black,fill=white,inner sep=2.4pt},
  active/.style={circle,draw=blue!65!black,fill=blue!25,inner sep=2.4pt},
  ghost/.style={diamond,draw=black,fill=black!15,inner sep=2.4pt}]
\draw[thick,black!55] (0,0)--(7,0);
\draw[very thick,blue!65!black] (1,0)--(2,0);
\foreach \j in {0,...,7}
  \node[font=\small,above] at (\j,.23) {$\j$};
\node[ghost] at (0,0) {};
\node[active] at (1,0) {};
\node[active] at (2,0) {};
\node[hit] at (3,0) {};
\node[circle,draw=black,dashed,fill=white,inner sep=2.4pt] at (4,0) {};
\node[active] at (5,0) {};
\node[hit] at (6,0) {};
\node[ghost] at (7,0) {};
\node[font=\scriptsize,align=center,below] at (0,-.22) {ghost\\birth};
\node[font=\scriptsize,below] at (1,-.22) {active};
\node[font=\scriptsize,below] at (2,-.22) {active};
\node[font=\scriptsize,align=center,below] at (3,-.22) {pending\\hit};
\node[font=\scriptsize,align=center,below] at (4,-.22) {pending\\miss};
\node[font=\scriptsize,below] at (5,-.22) {active};
\node[font=\scriptsize,align=center,below] at (6,-.22) {pending\\hit};
\node[font=\scriptsize,align=center,below] at (7,-.22) {ghost\\death};
\draw[decorate,decoration={brace,mirror,amplitude=4pt}]
  (.93,-.88)--(2.07,-.88)
  node[midway,below=5pt,font=\scriptsize] {first active run};
\node[font=\scriptsize,below] at (5,-.95) {second active run};
\end{tikzpicture}
\caption{A schematic lifespan of a prime from group $i$.
All active visits and both ghost endpoints have its common line;
the pending visit at $4$ misses that line. The line contributes
$1/(p+1)$ once. After extracting $b'_p$, the ghost signs,
pending-hit damping, and active-run normalizations give the weight
$\eta'_0\eta'_7q_3q_6 V_i^{-2}/((p+1)b'_p)$.
The two powers of $V_i^{-1}$ count active runs, not active visits.}
\label{fig:lifespan}
\end{figure}


\subsection{The symmetric memory space and its operations}

We realize the preceding expansion by storing each pending prime
together with its line. A later active run retrieves that stored pair,
instead of paying for a new line. The following measures and transfer
factors are chosen to reproduce \eqref{mem:run-charge}.

Let $\mathcal Z$ be the finite set of primitive integer vectors in the
tube which satisfy the fixed-root box condition; it has counting
measure. An active list is an ordered $M$-tuple in each group, endowed
with measure $\bigotimes_i\nu_i^{\otimes M}$. The underlying list space
allows repetitions. Distinctness within physical lists and the other
local restrictions will be imposed in the operators.

A memory particle in group $i$ consists of a prime and an anchor,
\begin{equation}\label{mem:particle-space}
 \mathcal X_i=\{(p,\ell):p\in\mathcal P_i,
                 \ \ell\in\mathbb P^1(\mathbb F_p)\},\qquad
 \lambda_i(p,\ell)=\nu_i(p).
\end{equation}
The line coordinate has counting measure, not uniform probability
measure. In particular, anchoring a particle at $[z]_p$ leaves the
prime measure $\nu_i(p)$, without another factor $1/(p+1)$. A particle
\emph{hits} $z$ when its anchor equals $[z]_p$.

For a vector of memory sizes $\boldsymbol k=(k_1,\ldots,k_K)$,
use symmetric functions of the $k_i$ particles in each group and the
measure $\lambda_i^{\otimes k_i}/k_i!$. The Hilbert space is
\begin{equation}\label{mem:hilbert-space}
 \mathscr H=
 \ell^2(\mathcal Z)\otimes
 L^2\!\left(\bigotimes_i\nu_i^{\otimes M}\right)
 \otimes\bigotimes_{i=1}^K
 \left(\bigoplus_{k_i\ge0}
 L^2_{\mathrm{sym}}\!\left(\mathcal X_i^{k_i},
                   \frac{\lambda_i^{\otimes k_i}}{k_i!}\right)\right).
\end{equation}
Products and sums of the operators below use the row convention:
the input state is fixed, choices leading to output states are summed,
and the resulting operator acts on a test function at the output.
Let $\mathbf b$ be one when $z=e_1$ and memory is empty, and zero
otherwise; it is independent of the active lists. By
\eqref{mem:measure-masses},
\begin{equation}\label{mem:boundary-norm}
 \norm{\mathbf b}^2=\prod_i s_i^M
   =\exp\!\left(O_K\!\left(\frac{M(N+1)}{p_{\min}}\right)\right)
   =1+o(1).
\end{equation}

Fix once and for all
\begin{equation}\label{mem:rho}
 0<\rho<1,\qquad \rho^2>\sqrt q.
\end{equation}
The ghost operation $\mathcal G_j$ leaves the position and active list
unchanged. Independently in each group it first chooses a subset of
the indexed pending particles which hit the current position and
deletes them, paying $-\eta'_j/\rho$ per deletion. Each surviving hit
pays $q_j/\rho^2$. It then appends an ordered batch of $l\ge0$ primes,
integrated with measure $\nu_i^{\otimes l}/l!$, anchored at the current
lines, and with coefficient $(-\eta'_jV_i/\rho)^l$.
The divisor $l!$ makes the batch an unordered birth set when the primes
are distinct; the ordered integral will also be useful after that
restriction is removed.

The edge $\mathcal E_j$ retains the geometric multiplier in
\eqref{eq:det-kernel}, or its adjoint according to the alternating
physical path, but removes its $\omega$-damping and its active
divisibility tests. It replaces target unshared counting and
normalization by the following operations between the same endpoint
slot symmetrizations.
\begin{enumerate}[label=(\roman*)]
 \item Pay $\rho$ for every old pending particle which hits the source
 $z$. Choose a subset of the unshared source slots to store in memory,
 anchored at $[z]_p$, and drop the other unshared source slots. Copy the
 shared $J$ slots in each group.
 \item Choose a target position $z'$ and its unshared label in each
 group. A subset of these slots is filled by promoting indexed old
 particles which hit $z'$, removing them from memory and paying
 $1/V_i$ per promotion in group $i$. Particles just stored on this edge
 cannot be promoted on the same edge. Every other unshared target slot
 is filled by an independent $\nu_i$ draw.
 \item Pay $\rho$ for every particle remaining in memory which hits
 $z'$, including newly stored particles. Impose the physical endpoint
 list restrictions, the cross-edge ban, the pad dyad, all good-state
 tests and geometric cutoffs, using the actual source and target label
 products.
\end{enumerate}
Positions are always summed with counting measure. Since the real
root matrix has determinant one, its edge determinant is
$\det(z,z')$.

For the moment impose also that all \emph{performed births} have
globally distinct prime values. Births comprise initial active
entries, fresh target entries, and ghost creations. This is a
restriction on the fully expanded history, and is not claimed to be
a local projection on $\mathscr H$. Then the independent-line path sum,
divided by its extracted baseline, is exactly
\begin{equation}\label{mem:operator-history}
 \inner{\mathbf b}
 {\mathcal G_0\mathcal E_0\mathcal G_1\cdots
  \mathcal E_{N-1}\mathcal G_N\mathbf b},
 \qquad\text{with global birth distinctness imposed.}
\end{equation}
Here equality refers to the choice-by-choice expansion of the right
side. To verify it, a shared prime continues with the same line:
$p\mid D$ implies $\det(z,z')=0\pmod p$, and primitivity implies
$[z]_p=[z']_p$. A prime leaving activity either terminates there or is
stored. If it is used again, global birth distinctness forces its
promotion from that stored particle. Ghost creations and deletions
are exactly the optional endpoints of
\eqref{mem:orphan-identity}--\eqref{mem:suffix}.

At an internal pending hit, the incoming edge, ghost survival, and
outgoing edge give
\begin{equation}\label{mem:rho-cancellation}
 \rho\,(q_j/\rho^2)\,\rho=q_j.
\end{equation}
At a ghost birth the creation coefficient and outgoing hit factor
give $-\eta'_jV_i$; at termination the incoming hit factor and deletion
give $-\eta'_j$. There are no unmatched factors at $0,N$, because memory
starts and ends empty. An active birth has measure
$\nu_i(p)=((p+1)b'_p)^{-1}V_i^{-1}$, a ghost birth has prime measure
$V_i\nu_i(p)=((p+1)b'_p)^{-1}$ after this cancellation, and each
promotion pays $V_i^{-1}$. These are exactly the factors in
\eqref{mem:run-charge}: each new active run contributes $V_i^{-1}$,
while the lifespan has only one line-probability factor.
Conversely a compatible primewise lifespan determines all these
choices: store on leaving any nonfinal active run, promote on entering
the next, and create or terminate at its specified ghost endpoints.
The factorial birth integral counts its unordered ghost birth set
once. This proves the identity, including primes which leave and
later re-enter activity.

\subsection{Truncation and adjoints}

Let $B=\lceil L^2\rceil$ and let $\Pi_B$ restrict total memory size to
at most $B$. We may insert $\Pi_B$ before and after every operation in
\eqref{mem:operator-history}, with error
\begin{equation}\label{mem:truncation-error}
 O_A(L^{-AN})\qquad\text{for every fixed }A>0.
\end{equation}
This assertion is made while births are still globally distinct.
There are at most $KN$ stores in a history. A history reaching memory
size greater than $B$ has therefore had at least $B-KN$ ghost births.
Insert a factor two at each ghost birth in an absolute majorant.
For a prime active somewhere, there are at most $(N+2)^2$ possible
ghost endpoint pairs. There are $O_K(M+N)$ such primes along a fixed
active path. For a never-active prime, fixing its earlier ghost endpoint
and summing its possible later endpoints costs $O_q(1)$: successive
later endpoints with the required line acquire successive internal
factors at most $\sqrt q<1$. Summing the earlier endpoint gives a
total $O_q((N+1)/p)$ for the absolute ghost-only contribution, including
the inserted factor two. Consequently the product of all never-active
prime costs is $\exp(O_K(N+1))$.

The extracted baseline and its reciprocal have logarithms
$O_K(N+1)$, by \eqref{mem:baseline} and the bounded reciprocal prime
sums. The path and label enumeration used in the root replacement
costs $\exp(O_K(L^{.73}))$; multiplying by the active endpoint choices
and the preceding absolute costs gives, with room to spare,
$\exp(O_K(L^{.74}))$. Removing the inserted factors on the omitted
histories gives the upper bound
\[
 2^{-B+KN}\exp(O_K(L^{.74})),
\]
which proves \eqref{mem:truncation-error}.

We next omit global birth distinctness and work on
$\mathscr H_B=\Pi_B\mathscr H$. The resulting signed norm estimates
apply only to the unrestricted history; we will restore distinctness
by an exact expansion in equality constraints. Repeated pending values
are now permitted. All deletion and promotion choices remain choices
of indices, so that the following
adjoint identities hold also in their presence.

Suppose $a$ particles are deleted from a memory list of size $k+a$.
For an unordered ghost deletion batch its indexed choices and the
factorial memory measure give
\begin{equation}\label{mem:ghost-factorial}
 \frac{\binom{k+a}{a}}{(k+a)!}=\frac{1}{k!a!}.
\end{equation}
Its anchors are forced by the hit condition, leaving exactly the
$\nu_i$ integrations. A fresh ghost batch has the same factorial
integral. Interchanging the deleted and appended batches therefore
gives the adjoint ghost operation: its deletion coefficient is
$-\eta'_jV_i/\rho$ and its creation coefficient is $-\eta'_j/\rho$;
the surviving-hit multiplier stays $q_j/\rho^2$.

For retrieval into $a$ prescribed ordered active slots the corresponding
identity is
\begin{equation}\label{mem:edge-factorial}
 \frac{\binom{k+a}{a}a!}{(k+a)!}=\frac{1}{k!}.
\end{equation}
The retrieved anchors again force one line and leave $\nu_i$ prime
integrations. Newly stored particles get those same prime measures
from their source active entries. Thus at two fixed positions the
joint measure of all source and target active entries and all memory
coordinates is unchanged on exchanging stores and promotions. The
factor $1/V_i$ on a forward promotion stays on that transfer, becoming
a reversed store factor; this is a bounded factor. The two endpoint
$\rho$ factors exchange roles. Summing over the two positions
preserves the equality because their measure is counting measure.
Slot symmetrizations are self-adjoint averages, and all endpoint and
edge restrictions transpose with the kernel. These observations
prove the adjoint rules even when list restrictions depend jointly
on both positions.

For later use, a local choice multiplier of modulus at most one may be
inserted provided it is equivariant under simultaneous permutations of
the old memory coordinates and the selected deletion or promotion
indices. The indexed choice sum then preserves symmetry in the old
particles. In a ghost birth integral, a multiplier depending on the
ordered fresh batch is averaged over permutations of that batch when
acting on symmetric test functions. This leaves the integral unchanged,
keeps its modulus at most one, and retains the factor $1/l!$. These
modified operations therefore act on the same symmetric spaces. In
both adjoint directions they are dominated in absolute value by the
positive kernels obtained by summing the absolute values of the
unmodified operation-choice coefficients. An arbitrary multiplier depending
on an old particle's index alone need not preserve symmetry and is not
allowed here.

\subsection{Absolute bounds for ghosts and edges}

Call an edge \emph{clean} if it stores and promotes no particle, and
\emph{dirty} otherwise. A clean edge draws all unshared target labels
freshly, and its signed minor-arc kernel will give cancellation.
For a dirty edge the good-state tests instead make an absolute row or
column sum small. We first bound the ghosts and record crude edge
bounds, then prove these two kinds of small edge estimate.

For a state with total pending size $k_{\mathrm{tot}}$ use the positive
Schur weight
\begin{equation}\label{mem:schur-weight}
 w=v_0^{k_{\mathrm{tot}}},
\end{equation}
where $v_0$ is a sufficiently large fixed constant. A weighted row
bound $R$ means that the absolute row sum against output weight is
at most $R$ times input weight; the analogous bound for the adjoint is
the weighted column bound $C$. Weighted Schur gives norm at most
$\sqrt{RC}$. Restrictions by $\Pi_B$ can only decrease these absolute
sums.

In one group with $h$ pending hits, the ghost row bound, divided by
the input weight, is
\begin{equation}\label{mem:ghost-row}
 \exp\!\left(\frac{v_0\eta'_jV_i s_i}{\rho}\right)
 \left(\frac{q_j}{\rho^2}+\frac{\eta'_j}{\rho v_0}\right)^h.
\end{equation}
The exponential sums fresh births; the power sums deletion or survival
of each old hit. The adjoint bound is
\begin{equation}\label{mem:ghost-column}
 \exp\!\left(\frac{v_0\eta'_j s_i}{\rho}\right)
 \left(\frac{q_j}{\rho^2}+\frac{\eta'_jV_i}{\rho v_0}\right)^h.
\end{equation}
Since $q_j\le\sqrt q<\rho^2$ and the $V_i$ are bounded above and below,
one fixed $v_0$ makes both parentheses at most one, for all $i,j$ and
large $x$. Multiplication over the fixed groups proves
\begin{equation}\label{mem:ghost-norm}
 \norm{\mathcal G_j}_{\mathscr H_B\to\mathscr H_B}\le C_K,
\end{equation}
with the same absolute row and column bounds under the bounded choice
modifications just described.

For the crude edge bounds fix the source $z$ and choose an integral
complement $w_z$ with $\det(z,w_z)=1$. Write
\begin{equation}\label{mem:neighbor}
 z'=j_1w_z+h z,\qquad j_1=\det(z,z'),\qquad
 \left|h+j_1\frac{\tau(w_z)}{\tau(z)}\right|\le16.
\end{equation}
The last inequality follows from $\tau(z')/\tau(z)$ lying in the
fixed box range. Enlarging its absolute constant would have no effect
on any argument below. For fixed source, labels, and determinant
$j_1$, this permits $O(1)$ targets. In the raw determinant part
$j_1=D(b-a)$ is fixed. In the comparison part $j_1=Dt$ has $O(Y)$
possibilities, and its multiplier is at most
\begin{equation}\label{mem:major-absolute}
 C\abs{\mathfrak M}\ll \frac{L^{3A_0}}{Y}.
\end{equation}
The masses of all fresh label draws are bounded by $C_K$.
There are at most $C_KB^K$ indexed promotion choices. Store subsets,
Schur weight ratios and $V_i$ factors cost $C_K$. Source and target
symmetrizations are probability averages. Using the adjoint
calculation for columns, we obtain a constant $C=C(K,A_0,q)$ such that
\begin{equation}\label{mem:edge-crude}
 R(\mathcal E_j),\ C(\mathcal E_j)
       \le C_K B^K(1+L^{3A_0})\le L^C.
\end{equation}
This holds for arbitrary multipliers of modulus at most one on the
choices. It also holds after fixing any or all fresh prime values,
with their measures left outside the bound. This uniformity will
preserve the cost of a forced prime atom in the birth-distinctness
argument.

\subsection{The lattice box and the clean signed norm}

We first extract a geometric consequence of the first good-state
test. Partition $\mathcal Z$ by intervals of length $H=d_0Y$ for
$z_2/\tau(z)$. The identity
\[
 \frac{z'_2}{\tau(z')}-\frac{z_2}{\tau(z)}
      =\frac{\det(z,z')}{\tau(z)\tau(z')}
\]
shows that an edge connects only cells whose indices differ by a
bounded amount. At fixed shared list, hence fixed squarefree product
$D$, partition further by the tuple $([z]_p)_{p\mid D}$. An edge
preserves this tuple. The resulting operator is a sum over a bounded
number of cell-index shifts of direct sums of blocks. It suffices to
bound each block uniformly.

Choose a relevant good source $z$ in one such pair of cells and a
complement $w_z$ as in \eqref{mem:neighbor}. All vectors in either cell
with the specified line tuple belong to the lattice
\[
 \Lambda=\mathbb Z z+\mathbb Z D w_z.
\]
Indeed, on writing a vector as $h z+jw_z$, equality of the projective
lines modulo every $p\mid D$ implies $p\mid j$; squarefreeness gives
$D\mid j$. Its coordinates $(h,l)$ in this lattice satisfy
\begin{equation}\label{mem:sheared-box}
 |l|\ll Y,\qquad |h+l\alpha|\ll1,
 \qquad \alpha=D\frac{\tau(w_z)}{\tau(z)}.
\end{equation}
The $l$ bound follows from the cell separation and $D\asymp d_0$;
the other follows from the box condition.

The image of this integer lattice under
$(h,l)\mapsto(h+l\alpha,l/Y)$ has covolume $1/Y$.
Let $\lambda$ be the length of its shortest nonzero vector. If
$\lambda<Y^{-.8}$, its second coordinate gives
$0<|l|<Y^{.2}$, and its first gives
$\|l\alpha\|<Y^{-.8}<Y^{-.7}$; $l=0$ is impossible for a nonzero
vector this short. This contradicts the first good-state test, since
$\tau(w_z)/\tau(z)$ is a lift of $r_{gz}$. Thus
$\lambda\ge Y^{-.8}$. Dirichlet's elementary pigeonhole approximation,
using denominators at most $\lceil\sqrt Y\rceil$, gives
$\lambda\ll Y^{-.5}$.

A shortest vector is primitive in the lattice, so complete it to a
basis. Subtract a multiple of the first vector from the second to
make its parallel component at most $\lambda/2$. Its perpendicular
component is $1/(Y\lambda)$; since $\lambda^2\ll1/Y$, its length is
$O(1/(Y\lambda))$. Inverting this basis on the bounded region in
\eqref{mem:sheared-box} puts all relevant vectors in a centered
coordinate box with side parameters $O(1/\lambda)$ and $O(Y\lambda)$.
Choose its orientation positively. Enlarging constants, the box has
parameters $A',B'$ satisfying
\begin{equation}\label{mem:lattice-shape}
 A'B'\ll Y,\qquad Y^{.1}\le A',B'\le Y^{.9}.
\end{equation}
The change of integer basis has determinant one, so that
$\det(z,z')/D$ is the ordinary determinant of their new integer
coordinates. This proof applies unchanged for a continuous root:
the tested quantity $r_{gz}$ is still
$\tau(w_z)/\tau(z)\bmod1$.

For a clean edge, before the slot symmetrizations, condition on the
shared list and all pending particles. They are unchanged by the central
operation. Its remaining
pending-hit factors are diagonal contractions at the two endpoints.
We may omit the extra cross-edge ban at a norm cost smaller than any
power of $L^{-1}$. In fact, keeping distinctness of each full active
list, a forbidden coincidence requires one of the fresh labels to
equal a fixed source label, costing at most
$C_KM\sup_{i,p}\nu_i(p)$. After the labels are fixed,
\eqref{mem:neighbor} and \eqref{mem:major-absolute} give an absolute
row bound $C_K L^{3A_0}$; the same argument for the adjoint gives the
column bound. Equation~\eqref{mem:measure-masses} proves the claimed
negligible norm cost.

Inside a lattice block, the unshared ordered list, containing one
prime from each disjoint band, is uniquely determined by its product
$b$. Its measure is
\begin{equation}\label{mem:product-mass}
 m(b)=\prod_{i=1}^K\nu_i(p_i)\le\frac{C_K}{Y}
 \quad\text{when }\eta(b/Y)\ne0.
\end{equation}
Conjugation from $L^2(m)$ to counting measure inserts
$\sqrt{m(b)m(a)}$. Consequently the clean norm is at most $C_K$ times
the norm of the following operator on the integer box in
\eqref{mem:lattice-shape} and an unrestricted integer product variable:
\begin{equation}\label{mem:clean-kernel}
 \frac1Y\psi\!\left(\frac{\det(\boldsymbol h,\boldsymbol l)}Y\right)
 \int_{\mathbb T\setminus\mathfrak M}
 \e\!\left(\theta
      (\det(\boldsymbol h,\boldsymbol l)-b+a)\right)\,\dd\theta.
\end{equation}
Here $\mathbb T=\mathbb R/\mathbb Z$. Extending the product variable
means composing with extension by zero and restriction. All actual
product supports, full-list restrictions, goodness tests, cutoff
factors, and the bounded square-root factors from
\eqref{mem:product-mass} are endpoint diagonal multipliers. The factor
$d_0/D$ is bounded. Thus none of these operations increases the bound
except by $C_K$. We use the transposed conjugate kernel on a reversed
edge.

Fourier transformation in the unrestricted integer product variable
diagonalizes its convolution: for each
$\theta\in\mathbb T\setminus\mathfrak M$ the fiber is the determinant
exponential matrix with its smooth cutoff, divided by $Y$. If
$\widehat\psi(u)=\int_{\mathbb R}\psi(t)\e(-ut)\,\dd t$, then
\begin{equation}\label{mem:cutoff-fourier}
 \psi\!\left(\frac{\det(\boldsymbol h,\boldsymbol l)}Y\right)
 \e\bigl(\theta\det(\boldsymbol h,\boldsymbol l)\bigr)
 =\int_{\mathbb R}\widehat\psi(u)
   \e\bigl((\theta+u/Y)\det(\boldsymbol h,\boldsymbol l)\bigr)
   \,\dd u.
\end{equation}
The density $\widehat\psi$ decreases faster than any inverse power.

For $|u|\le L^{A_0}/2$, put $\vartheta=\theta+u/Y$. Dirichlet
approximation with $Q=\lceil Y/L^{A_0}\rceil$ gives a reduced
$c/d$ such that
\begin{equation}\label{mem:dirichlet-minor}
 L^{A_0}<d\le Q,\qquad
 \left|\vartheta-\frac cd\right|
        \le\frac1{dQ}\le\frac1{d^2}.
\end{equation}
Indeed $d\le L^{A_0}$ would put $\theta$ within
$\tfrac32L^{A_0}/Y$ of a rational of allowed denominator, contrary
to the definition of $\mathfrak M$ with radius $2L^{A_0}/Y$.

For completeness, let $T_{A',B'}(\vartheta)$ have entries
$\e(\vartheta hk)$ for integer $|h|\ll A'$, $|k|\ll B'$.
Its Gram matrix and the geometric-sum estimate give
\begin{equation}\label{mem:bilinear-gram}
 \norm{T_{A',B'}(\vartheta)}^2
 \ll\sum_{|h|\ll A'}\min\!\left(B',\frac1{\|h\vartheta\|}\right)
 \ll (A'/d+1)\bigl(B'+d\log(2d)\bigr).
\end{equation}
In the term $h=0$ the minimum is interpreted as $B'$. To justify the
last inequality, split the $h$ interval into blocks of length at most
$d/2$. For two different indices in a block, reduction of $c/d$ and
the error in \eqref{mem:dirichlet-minor} separate their fractional
parts by at least $1/(2d)$. Ordering their distances to the nearest
integer bounds each block by
$O(B'+d\sum_{1\le j\le d}j^{-1})$.

Up to a permutation of columns, the determinant matrix
$\e(\vartheta(h_1l_2-h_2l_1))$ is the tensor product of
$T_{A',B'}(\vartheta)$ and its transposed conjugate. Its norm is
therefore the squared norm in \eqref{mem:bilinear-gram}. After division
by $Y$, \eqref{mem:lattice-shape} implies the bound
\begin{equation}\label{mem:minor-decay}
 \frac{\norm{(\e(\vartheta\det(\boldsymbol h,\boldsymbol l)))}}Y
 \ll \frac1d+Y^{-.1}\log(2Y)
          +\frac{d\log(2d)}Y
 \ll_K L^{.2-A_0}+Y^{-.1}\log(2Y).
\end{equation}
For the complementary Fourier tail in \eqref{mem:cutoff-fourier},
the trivial matrix norm is $O(A'B')=O(Y)$, so rapid decay of
$\widehat\psi$ gives any required negative power of $L$. Combining
the direct sums of cell blocks and the endpoint symmetrizations,
we conclude that, for any fixed $G>0$, choosing $A_0$ sufficiently
large in terms of $G$ gives
\begin{equation}\label{mem:clean-bound}
 \norm{\mathcal E_j^{\mathrm{clean}}}\le L^{-G}
\end{equation}
for large $x$, with an arbitrarily fixed margin in the exponent.
Constants depending on the later fixed value of $K$ are absorbed by
that margin and the threshold for $x$.

\subsection{Dirty raw edges: the two Schur sides}

Fix the set $S$ of stored groups and the set $I$ of promoted groups.
There are at most $4^K$ choices, a constant. All store and promotion
coefficients and the ratio of Schur weights cost $C_K$.

If $I=\varnothing$, the weighted raw row bound is $C_K$, independently
of $B$. For each source omission tuple, sample the fresh target label
in every group and use \eqref{mem:neighbor} with its fixed determinant.
There are no pending-index choices. The source symmetrization averages
the omission tuples and thus does not multiply their number. This
argument allows arbitrary $S$.

Suppose $I\ne\varnothing$ and put $i_* =\max I$. Fix the source state
and sample the fresh target labels in the groups outside $I$, with
product $T_f$. Their product law $\nu$ is boundedly comparable to the
law $\mu$ in the second good-state test, by
\eqref{mem:measure-masses}. Outside an exceptional mass
$C_K\exp(-\tfrac14L^{a_{i_*}})$, that test supplies its asserted phase
separation. Fix one old index for promotion in group $i_*$; there are
at most $B$ such indices, and write its prime and anchor as
$(p_*,\ell_*)$. Let $Z$ be the product of the other promoted numeric
primes and let $P_{\mathrm{full}}$ be the product of the complete source
active list. The raw determinant equation is
\begin{equation}\label{mem:dirty-determinant}
 j_1=P_{\mathrm{full}}-Dp_*ZT_f.
\end{equation}
The ban excludes $p_*$ from the entire source list. Thus
$j_1\equiv P_{\mathrm{full}}\ne0\pmod{p_*}$, independently of the
omission tuple and of $Z$.

In \eqref{mem:neighbor} the line condition
$[j_1w_z+hz]_{p_*}=\ell_*$ is either impossible or fixes a single
residue $h=h_*\pmod{p_*}$. To see uniqueness, in the basis
$(z,w_z)$ the second coordinate $j_1$ is nonzero modulo $p_*$;
therefore its projective line determines the first coordinate
uniquely. This residue depends on the fixed index and
$P_{\mathrm{full}}$, but not on $D,Z$. Substitution in
\eqref{mem:neighbor}, using the real lift
$r_z^*=\tau(w_z)/\tau(z)$, yields
\begin{equation}\label{mem:anchor-congruence}
 \left\|DZ T_f r_z^*
       -\frac{P_{\mathrm{full}}r_z^*+h_*}{p_*}\right\|
       \le\frac{16}{p_*}.
\end{equation}
The phase separation is $100e^{-L^{a_{i_*}}}$, whereas the diameter
of the interval in \eqref{mem:anchor-congruence} is at most
$32/p_*\le32e^{-L^{a_{i_*}}}$. Hence at most one distinct integer
$DZ$ is possible.

This also determines $D$ and $Z$ separately. All prime factors of $D$
are source labels and all prime factors of $Z$ are excluded from the
source list by the ban. Their prime factorizations therefore separate
unambiguously. Because the source list is distinct, $D$ determines
which single label was omitted in every group. Under source
symmetrization this one numeric omission tuple has mass exactly
$M^{-K}$. It is an average over tuples, including over their internal
orders, rather than an unnormalized sum. The disjoint prime bands
also make $Z$ determine each other promoted numeric prime.

The labels now fix $j_1$, so there are $O(1)$ target positions. At any
one of them, multiplicities among the remaining promoted indices are
absorbed by the outgoing pending damping. If a group has $h$ old
particles hitting that target, the choice of one eligible index and
the damping of all other old hits cost at most
\begin{equation}\label{mem:index-damping}
 h\rho^{h-1}\le\sup_{n\ge1}n\rho^{n-1}<\infty.
\end{equation}
This remains true for repeated numeric values and repeated anchors.
Additional newly stored hits only decrease the factor. Applying this
in each of the other promoted groups costs $C_K$. For the exceptional
fresh draws use the crude $B^K$ count and
\eqref{mem:neighbor}. We obtain the weighted raw row bound
\begin{equation}\label{mem:epsilon}
 \varepsilon_K=C_K\left(\frac{B}{M^K}
                         +B^K e^{-cL^{.1}}\right)
 \qquad(I\ne\varnothing).
\end{equation}

Reversal exchanges $S$ and $I$ by \eqref{mem:edge-factorial}, changes
only bounded $V_i$ factors, and retains both good-state tests and the
ban. Its raw determinant equation, with the reversed sign, again has
source full product minus $D$ times the target unshared product.
The preceding argument therefore gives the following complete table.
Every entry includes the fixed $C_K$ factors.
\begin{center}
\begin{tabular}{cccc}
\toprule
Stored groups $S$ & Promoted groups $I$ & Weighted row & Weighted column\\
\midrule
$\varnothing$ & $\varnothing$ & $C_K$ & $C_K$\\
$\varnothing$ & nonempty & $\varepsilon_K$ & $C_K$\\
nonempty & $\varnothing$ & $C_K$ & $\varepsilon_K$\\
nonempty & nonempty & $\varepsilon_K$ & $\varepsilon_K$\\
\bottomrule
\end{tabular}
\end{center}
The first row is only an absolute estimate; its signed estimate is
\eqref{mem:clean-bound}. Once $K$ is large enough that
$\varepsilon_K\le1$, weighted Schur bounds the sum of the three dirty
raw cases by
\begin{equation}\label{mem:raw-dirty-norm}
 C_K\sqrt{\varepsilon_K}
 \le C_K L^{1-.005K+o(1)}+O_A(L^{-A})
 \quad\text{for every fixed }A.
\end{equation}
The bounded entry opposite a small Schur entry is essential here: it
contains no factor $B^K$. Thus making $.005K$ larger than a prescribed
constant purchases that constant in log decay. All remaining
$K$-dependence in this estimate is a fixed multiplicative constant.

\subsection{Dirty comparison edges}

Use the supremum \eqref{mem:major-absolute}, discarding its dependence
on the target unshared labels. Fix a source omission tuple, hence
$D$, and suppose $I\ne\varnothing$. A promoted old index with prime
$p$ requires one specified projective line at the target. By the ban,
$p\nmid D$. The lattice basis used to obtain
\eqref{mem:lattice-shape} has determinant $D$ in the original integer
coordinates and is therefore invertible modulo $p$. The prescribed
line is consequently one nonzero homogeneous linear congruence in
the two box coordinates.

For a box of sides $A',B'$, the number of solutions of such a
congruence is
\begin{equation}\label{mem:line-count}
 O\!\left(\frac{A'B'}p+\max(A',B')\right)
 \ll Y\left(\frac1p+\frac1{\min(A',B')}\right).
\end{equation}
If the coefficient of the shorter coordinate is nonzero, fix the
other coordinate and count its solutions in that shorter interval.
If that coefficient vanishes, the congruence restricts the longer
coordinate instead; the same displayed upper bound follows. We have
included vectors divisible by $p$ in this count, which only increases
it.

For each source only a bounded number of neighboring cell blocks
occurs. Take the union over at most $B$ old indices in one promoted
group. The other index choices can be bounded by $B^{K-1}$, or by
\eqref{mem:index-damping}. Sum the fresh prime masses and the source
omission average. Equations \eqref{mem:major-absolute},
\eqref{mem:lattice-shape}, and \eqref{mem:line-count} give a row bound
\begin{equation}\label{mem:comparison-small}
 \zeta_K
   :=C_K L^{3A_0}B^K\left(p_{\min}^{-1}+Y^{-.1}\right),
 \qquad \zeta_K=O_A(L^{-A})
 \quad\text{for every fixed }A.
\end{equation}
All parameters $K,A_0$ are fixed in the last assertion. The opposite
Schur side is at most $L^C$ by \eqref{mem:edge-crude}. For a store-only
piece apply the same argument to the adjoint; with both stores and
promotions it applies to both sides. For clarity the comparison
counterpart of the dirty table is
\begin{center}
\begin{tabular}{cccc}
\toprule
Stored groups $S$ & Promoted groups $I$ & Weighted row & Weighted column\\
\midrule
$\varnothing$ & nonempty & $\zeta_K$ & $L^C$\\
nonempty & $\varnothing$ & $L^C$ & $\zeta_K$\\
nonempty & nonempty & $\zeta_K$ & $\zeta_K$\\
\bottomrule
\end{tabular}
\end{center}
Weighted Schur makes every dirty comparison piece smaller than every
fixed negative power of $L$.

Choose $G>E_0+3$. First choose $A_0$ so that the clean argument gives
more than $G$ powers of decay, and then choose $K$ so that
$.005K>G+2$. Combining the clean, dirty raw, and dirty comparison
estimates, with their fixed finite sums, proves
\begin{equation}\label{mem:signed-edge}
 \norm{\mathcal E_j}_{\mathscr H_B\to\mathscr H_B}\le L^{-G}
\end{equation}
before global birth distinctness is restored. If desired all strict
inequalities here can be enlarged by one to absorb the constants.
The crude exponent $C(K,A_0,q)$ in \eqref{mem:edge-crude} has not
entered the choice needed in \eqref{mem:raw-dirty-norm}.

\subsection{Restoring global birth distinctness}

The contractions just proved did not require globally distinct births.
We now restore that condition exactly. A union bound over one repeated
prime would not suffice against the absolute cost of a long path;
we instead separate equality constraints by their rank.

Allocate a finite ordered set $\mathcal B$ of potential birth
addresses as follows: the $KM$ initial slots; the $K$ canonical target
slots at each edge, before its final symmetrization; and $B$ ordered
fresh-batch indices per group at each ghost operation. The truncation
ensures that a ghost batch has length at most $B$. Each address has
a local \emph{performed flag}: at an edge it is performed exactly if
that slot is filled freshly, and in a ghost batch exactly if its index
does not exceed the batch length. Initial addresses are always
performed. Thus
\begin{equation}\label{mem:address-count}
 Q_{\mathrm{birth}}:=\abs{\mathcal B}
       \le KM+KN+K(N+1)B=L^{O(1)}.
\end{equation}
Order these addresses chronologically, with a fixed order within
each operation.

Let $\mathcal C$ be a collection of disjoint nonsingleton subsets of
$\mathcal B$. For each block $C\in\mathcal C$ require that every
address is performed and that their prime values are equal. Denote
this event by $E_C$. Then the exact distinctness indicator is
\begin{equation}\label{mem:partition-identity}
 \1_{\text{distinct performed birth values}}
 =\sum_{\mathcal C}
       \prod_{C\in\mathcal C}(-1)^{|C|-1}(|C|-1)!
       \prod_{C\in\mathcal C}\1_{E_C}.
\end{equation}
One proof expands the right side as the sum of signs of all
permutations of the performed addresses which preserve their values:
a nontrivial cycle on $C$ has sign $(-1)^{|C|-1}$ and there are
$(|C|-1)!$ such cycles. In each equal-value class of size at least two
the signs of permutations sum to zero, and for a singleton they sum
to one. Unperformed addresses cannot occur in a nontrivial cycle,
which is exactly enforced by the flags.

Define the rank
\begin{equation}\label{mem:rank-definition}
 r(\mathcal C)=\sum_{C\in\mathcal C}(|C|-1).
\end{equation}
A rank-$r$ collection involves at most $2r$ addresses. Its total
absolute coefficient mass, summed over all collections of that rank,
is at most
\begin{equation}\label{mem:rank-coefficients}
 Q_{\mathrm{birth}}^{2r}=L^{O(r)}.
\end{equation}
Indeed the coefficients count permutations with that cycle rank,
and each such permutation has a fixed canonical decomposition into
$r$ transpositions. Its ordered list of transpositions, drawn from
at most $Q_{\mathrm{birth}}^2$ possibilities, determines the
permutation.

We first bound a fixed system of equality constraints absolutely.
In each block choose its earliest address as pivot, and condition
chronologically on the entire earlier history and all known pivot
values. Each nonpivot performed birth must draw one prescribed prime.
By \eqref{mem:measure-masses} this costs an atom at most
\begin{equation}\label{mem:atom}
 \Delta:=\sup_{i,p}\nu_i(p)\le C_K e^{-L^{.1}}.
\end{equation}
If its prescribed value belongs to a wrong group, the term is zero.
No later endpoint conditioning is imposed in this absolute bound:
we drop the final return requirement and use weighted row iteration.
Thus dependence of intermediate positions on the earlier pivot
values cannot change the atom cost.

Here is the precise uniform estimate justifying that iteration.
For an edge with $s$ prescribed nonpivot fresh values, fix all fresh
values first. The possible target positions still number $O(1)$ in
the raw part or $O(Y)$ in the comparison part. The latter count is
compensated by \eqref{mem:major-absolute}. Consequently its weighted
absolute row is at most
\begin{equation}\label{mem:edge-atoms}
 C_K B^K(1+L^{3A_0})\Delta^s,
\end{equation}
uniformly in the input state and all fixed values. The remaining fresh
values have bounded total mass, and all local restrictions may be
discarded for this upper bound.

For a ghost in group $i$, set
$c_i=v_0\eta'_jV_i/\rho$. If $s$ specified ordered coordinates of
its new batch are prescribed, its birth row sum is at most
\begin{equation}\label{mem:ghost-atoms}
 \sum_{l\ge s}\frac{c_i^l\Delta^s s_i^{l-s}}{l!}
 \le(c_i\Delta)^s\exp(c_i s_i).
\end{equation}
The actual requirement that the batch reaches its largest named
index can only reduce this sum. Within a batch the earlier pivot
coordinates are integrated before the prescribed later ones, which
gives the same estimate. Deletion and survival have weighted factor
at most one by \eqref{mem:ghost-row}. Anchors add no factor $p+1$,
as their measure is counting measure. Initial slots have the same
single-atom estimate. Thus for each fixed rank-$r$ constraint system,
the absolute boundary contribution is at most
\begin{equation}\label{mem:rank-absolute}
 L^{C(N+1)}(C_K e^{-L^{.1}})^r,
\end{equation}
for a fixed $C=C(K,A_0,q)$. The boundary weight is one because memory
is empty; after dropping return, its terminal indicator is at most
the positive Schur weight, so the chronological row bounds apply.

Choose, after $C$ is fixed,
\begin{equation}\label{mem:rank-cutoff}
 r_0=\left\lceil C'\frac{N\log L}{L^{.1}}\right\rceil,
\end{equation}
with $C'$ sufficiently large. Combining
\eqref{mem:rank-coefficients} and \eqref{mem:rank-absolute}, and using
$\log Q_{\mathrm{birth}}=O(\log L)$, bounds all ranks $r\ge r_0$ by
\begin{equation}\label{mem:high-rank}
 L^{C(N+1)}\sum_{r\ge r_0}
   \bigl(C_K Q_{\mathrm{birth}}^2e^{-L^{.1}}\bigr)^r
 \le L^{-(E_0+2)N}.
\end{equation}
The series is geometric for large $x$. For example its ratio is at
most $e^{-L^{.1}/2}$, so choosing $C'/2>C+E_0+3$ suffices after
enlarging the threshold for $x$.

For ranks $r<r_0$, use signed contraction on most edges. If
$C=\{a_0,a_1,\ldots,a_s\}$ is one block, with earliest address $a_0$,
write its equality event by ordinary circle orthogonality as
\begin{equation}\label{mem:local-phases}
 \1_{E_C}
 =\left(\prod_{a\in C}\1_{a\ \mathrm{performed}}\right)
   \int_{\mathbb T^s}
      \prod_{h=1}^s
          \e\bigl(\theta_h(p_{a_h}-p_{a_0})\bigr)
      \,\dd\theta_1\cdots\dd\theta_s.
\end{equation}
When an address is unperformed the integrand is interpreted as zero;
no value is assigned to an unperformed birth. After the circle
variables are fixed, every phase is attached only to the operation
where its prime was created. In particular the pivot receives the
single local multiplier
$\e(-p_{a_0}\sum_h\theta_h)$. Its value need not be remembered through
later operations. Flags are local choice restrictions of modulus at
most one. An edge flag tests whether its slot is freshly drawn rather
than promoted; it does not distinguish old promotion indices. A ghost
flag tests whether its fresh batch reaches the specified address.
Thus these flags and birth-value phases satisfy the old-memory
equivariance condition stated after \eqref{mem:edge-factorial}.
For fixed circle variables, backward induction through the later
operations gives a symmetric continuation function. Averaging the
ghost-batch multiplier over its new coordinates consequently leaves
the full integral unchanged, without increasing its absolute kernel,
changing its factor $1/l!$, or modifying any later operation.

It follows that a rank-$r$ system modifies at most $2r$ operations,
and hence at most $2r$ edges. Initial phases change $\mathbf b$ by a
bounded multiplier and preserve \eqref{mem:boundary-norm}. Every
ghost, modified or not, has norm at most $C_K$; a modified edge has
norm at most $L^C$ by \eqref{mem:edge-crude}; all other edges retain
\eqref{mem:signed-edge}. For each fixed set of phases the boundary
matrix element is therefore at most
\begin{equation}\label{mem:low-rank-term}
 C_K^{N+1}\norm{\mathbf b}^2
       L^{-G(N-2r)} L^{2Cr}.
\end{equation}
Integration of phases does not increase this bound. As
$r_0/N=O(\log L/L^{.1})=o(1)$, summing
\eqref{mem:low-rank-term} with \eqref{mem:rank-coefficients} gives
\begin{equation}\label{mem:low-rank}
 \sum_{r<r_0}Q_{\mathrm{birth}}^{2r}
 C_K^{N+1}\norm{\mathbf b}^2 L^{-G(N-2r)+2Cr}
       \le L^{(-G+o(1))N}.
\end{equation}
This remains true even though $C$ depends on $K$: $K$ was fixed
before letting $x$ grow, and its crude exponent is paid on only
$o(N)$ edges.

Equations \eqref{mem:high-rank} and \eqref{mem:low-rank} bound
\eqref{mem:operator-history} with exact global birth distinctness,
uniformly in all archimedean root variables. Multiply back the
baseline, whose modulus is at most one. The root integral has bounded
mass. The error \eqref{mem:truncation-error} and the root-replacement
error are $O_A(L^{-AN})$ for arbitrary fixed $A$. Because
$G>E_0+3$, these margins absorb all constants and prove
\eqref{mem:moment-bound}. The passage from this moment to
\eqref{eq:det-pairing} was proved in the preceding section.

\subsection{Removing goodness and stripping the pads}

We complete the second assertion of
Proposition~\ref{prop:minor-moment}. First remove the good-state
projections from the endpoint pairing. The part where the source is
bad, and separately the part where the target is bad, may be
majorized absolutely; bound the endpoint coefficients by their
log-power suprema. Apply the single-edge version of the root
replacement. In the independent-line model, a shared prime requires
the same line at both positions, or contributes zero. Every distinct
prime in the union of the two active lists is charged exactly once
by $1/(p+1)$. The ban makes the union consist of $M+1$ distinct primes
per group. Its list normalization is $V_i^{-(M+1)}$, so summing these
probabilities over ordered union lists is at most
\begin{equation}\label{mem:one-edge-label-mass}
 \prod_{i=1}^K
 \left(\sum_{p\in\mathcal P_i}\frac1{V_i(p+1)}\right)^{M+1}
 \le1.
\end{equation}
We have dropped any incompatibilities for an upper bound. Damping
is also at most one on the physical states and may be discarded.

For fixed labels, \eqref{mem:neighbor} counts $O(1)$ raw targets;
the comparison target count is $O(Y)$ and is compensated by
\eqref{mem:major-absolute}. For every fixed source list, the set of
root ratios failing either good-state test has measure
$O(e^{-cL^{.1}})$ by Lemma~\ref{det:good-state}. This estimate
is uniform in the other root coordinates. The preceding target
counts hold for every ratio, so they may be integrated over that bad
set. Reversal proves the identical assertion for failure of target
goodness. Equation~\eqref{mem:one-edge-label-mass} then gives the
normalized pairing error
\begin{equation}\label{mem:goodness-error}
 O\!\left(UV L^{C_1+3A_0}e^{-cL^{.1}}\right)
       +O_A(UV L^{-A}),
\end{equation}
where $C_1$ comes from the endpoint suprema. The root-replacement
error can here be made smaller than every log power by its
single-edge estimate; this also handles the indicators of goodness
failure. Therefore removing the projections costs
$O_A(UV L^{-A})$ for every fixed $A$.

The endpoint vectors are invariant under permutations of their
lists. Hence the two slot symmetrizations disappear when the pairing
is expanded. Write the ordered pads as $p_{i,1},\ldots,p_{i,J}$ in
group $i$, with total product $D$. The remaining source and target
slots have products $b$ and $a$. The common state normalization and
the fresh-target normalization are
$\prod_i V_i^{-(J+2)}$. Divide the pairing for the dyad $d_0$ by
$d_0$. Its scalar multiplier becomes $1/D$, and the exact measure
identity is
\begin{equation}\label{mem:pad-measure}
 \frac1D\prod_{i=1}^K V_i^{-(J+2)}
   =\left(\prod_{i=1}^K\prod_{h=1}^J\mu_i(p_{i,h})\right)
      \prod_{i=1}^K V_i^{-2}.
\end{equation}
Thus the pads are independent ordered harmonic draws, and the
remaining factors are exactly the two unshared-list normalizations
in the expanded square. There is no factorial or $M$-dependent
normalization left over: the pads occupy prescribed ordered shared
slots and the unshared label occupies its prescribed remaining slot.

The change of variables is $P=(Dbm,s)$, $Q=(Dar,n)$, so
$\det(P,Q)/D=bmn-ars$. The support and roughness properties of the
endpoint vectors give the original coefficients and make the
physical damping equal to one. Primitivity and the box conditions
are automatic on their nonzero support, as established in the
preceding section. On summing the ordinary pad dyads every ordered
pad tuple occurs once. Therefore, apart from the pad collision
restrictions, \eqref{mem:pad-measure} gives exactly the signed
difference between the determinant condition and
$\mathcal H_{\mathfrak M}$ in the expanded square.

For fixed disjoint unshared products $a,b$, independent pad draws
violate pad distinctness or the ban with probability at most
\begin{equation}\label{mem:pad-collisions}
 C_K M^2\sup_{i,p}\mu_i(p)
       \ll_K M^2e^{-L^{.1}}.
\end{equation}
This follows by a union bound over pairs of pad addresses and over
coincidences with the two fixed unshared primes in each group; for
two independent draws their equality probability is at most the
largest atom. To remove this probability loss from a signed sum,
we need absolute bounds for each of its two parts.

The absolute raw expanded square is
\begin{equation}\label{mem:raw-absolute-square}
 \ll_K XY L^{C_2},
\end{equation}
with $C_2$ depending on the coefficient bounds and fixed divisor
moments, independently of $K$. Indeed for each pair $a,b$ its original
$h$ representation has $h\ll X/Y$ and contributes at most a
log-power coefficient bound times
$\sum_{h\ll X/Y}\tau(1+ah)\tau(1+bh)$.
Cauchy and the divisor moment estimate from the preceding section
bound this by $(X/Y)L^{C_2}$. There are $O(Y^2)$ possible numeric
products $a,b$, and their normalization costs only $C_K$.

For the absolute comparison sum, fix $a,b\asymp Y$. If $k=mn$ and
$l=rs$, the cutoff imposes
\[
 |bk-al|\ll Y,\qquad k,l\ll X.
\]
For each $k$ there are $O(1)$ choices of $l$, and conversely, because
$a,b\asymp Y$. Thus Cauchy on this bounded-degree relation and the
ordinary divisor second moment give
\begin{equation}\label{mem:comparison-representations}
 \sum_{\substack{k,l\ll X\\|bk-al|\ll Y}}\tau(k)\tau(l)
 \ll\sum_{k\ll X}\tau(k)^2\ll X L^{C_3}.
\end{equation}
The endpoint coefficients cost a fixed log power, and
\eqref{mem:major-absolute} contributes $O(L^{3A_0}/Y)$. Summing the
$O(Y^2)$ possible $a,b$ proves
\begin{equation}\label{mem:comparison-absolute-square}
 \ll_K XY L^{C_4+3A_0},
\end{equation}
where $C_3,C_4$ are again independent of $K$. Multiplying
\eqref{mem:raw-absolute-square} and
\eqref{mem:comparison-absolute-square} by
\eqref{mem:pad-collisions} gives $O_A(XY L^{-A})$ for every fixed
$A$. This strips all pads.

The comparison estimate also allows the two unshared products to have
a common prime.
The number of pairs $a,b\asymp Y$ sharing any group prime is at most
\begin{equation}\label{mem:unshared-collisions}
 C_KY^2\sum_{p\in\bigcup_i\mathcal P_i}\frac1{p^2}
       \ll_K Y^2e^{-L^{.1}}.
\end{equation}
For fixed $p$ there are at most $O(Y/p)$ possible numeric products
divisible by $p$, and each product determines its prime list.
Apply \eqref{mem:comparison-representations} to each such pair and
then \eqref{mem:major-absolute}; their total is again smaller than
$XY$ times every fixed negative log power.

Finally, if $C_5$ denotes the log-power loss in
\eqref{eq:det-pairing}, its endpoint calculation gives $C_5$ depending
only on the fixed coefficient bounds, independently of $K$.
Since $UV=d_0YX$, division by $d_0$ bounds each pad dyad by
$XY L^{-E_0+C_5}$. There are $O_K(L)$ pad dyads, so their sum is
\begin{equation}\label{mem:final-replacement}
 O_K(XY L^{-E_0+C_5+1})+O_A(XY L^{-A}).
\end{equation}
Given a desired replacement precision $A$, choose
$E_0>A+C_5+2$ first, then choose $A_0$, then $K$ as above.
All exponentially small errors have already been bounded after these
parameters are fixed and impose no further condition on $E_0$.
This proves the replacement assertion of
Proposition~\ref{prop:minor-moment} and leaves the unrestricted major
term for the next section.
